%---------------------------------------------------------------%
%----------------- ELEMENTOS PRÉ-TEXTUAIS ----------------------%
%---------------------------------------------------------------%

% Capa oficial configurada para o Projeto Físico
\imprimircapa

%---------------------- Resumo ---------------------------------%
\begin{resumo}
Ferramentas de análise estática de segurança (SAST), como o Semgrep, apontam vulnerabilidades no código-fonte, mas a correção continua manual. Produtos comerciais já geram correções com modelos de linguagem (LLMs), e estudos recentes mostram que muitas dessas correções compilam e preservam o comportamento do sistema, mas não eliminam a falha. Este projeto especifica um agente autônomo de remediação que recebe o relatório do SAST e, para cada achado, investiga o repositório com ferramentas de busca e leitura, consulta uma base de conhecimento especializada por classe de vulnerabilidade e uma memória das correções já aceitas no projeto, propõe um \textit{patch} compatível com a arquitetura existente e o submete a uma verificação em camadas: compilação, testes existentes, reanálise por dois analisadores, teste de exploração que falha antes e passa depois da correção, preservação de comportamento e restrições de arquitetura. O resultado é entregue com o código, o raciocínio e um laudo de verificação, e só é aplicado após aprovação humana. O contexto do projeto é mantido por um índice atualizado incrementalmente. Nesta fase, o trabalho se concentra em projetos C\#/.NET e Java. Como não há um conjunto público de vulnerabilidades em C\# com testes de exploração, o trabalho também constrói esse conjunto; em Java, usa o Vul4J. A avaliação compara três configurações: LLM com contexto mínimo, LLM com contexto recuperado e agente completo.

\textbf{Palavras-chave}: Remediação de vulnerabilidades. Agentes de LLM. Análise estática de código. Verificação de \textit{patches}. Contexto de projeto.
\end{resumo}

%---------------------- Sumário --------------------------------%
\pdfbookmark[0]{\contentsname}{toc}
\enlargethispage{2\baselineskip}
\tableofcontents*
\cleardoublepage
